\section{Governed Execution}
\label{sec:governance}

\subsection{Identity before action}
Target resolution answers where to act; identity verification asks whether the
resolved target belongs to the intended entity. Structured text is authoritative
where available. On pure-pixel substrates, OCR ambiguity can make distinct
identifiers observationally identical. The safe response is abstention and
halt, not a guessed click. Importantly, identity guarantees apply only to armed
steps. Current real bundles can contain unarmed clicks, so coverage is reported
per workflow and certification policy can impose a minimum.

\subsection{Effect verification}
A screen may paint success after a partial, duplicate, rejected, or stale write.
Consequential steps can therefore declare typed effects and use an independent
verifier against REST, FHIR, or document state. The verifier snapshots the
system of record before and after an action and returns confirmed, refuted, or
indeterminate. A declared effect without a configured verifier halts.

Crucially, the strong-oracle read is performed \emph{out of band}: the verifier
queries the system of record through its own interface (a SQL read on a
read-only role, a REST or FHIR read, or a document or file read), independently
of the substrate that carried the write. It does not parse the pixel stream or
scrape the rendered page. The same effect contract therefore holds whether the
write was driven through a browser DOM, native accessibility, or a pixel-only
remote session: what changes is only the connector, not the check. Because the
read is out of band, an unreachable, unreadable, or unauthenticated system of
record maps to indeterminate and halts; the runtime never upgrades ``I could not
read the record'' into a reported success.

Two consequences follow for pixel-only substrates such as RDP and Citrix. First,
the effect check still applies as long as some out-of-band read path exists: the
RDP qualification in Section~\ref{sec:results}, for instance, reads the persisted
file through guest tools, a channel independent of the RDP display. Second, when
a locked-down session exposes \emph{no} such path, no database, no reachable API,
and no readable file, the transactional strong oracle is simply unavailable. The
only fallback is a same-application re-read of the record, which cannot catch a
partial save, a duplicate, a stale overwrite, or a lost update, and the honest
runtime treats that as a halt rather than a confirmation. On such substrates the
safety guarantee rests instead on the identity gate and halt-on-ambiguity, not on
effect verification. Effects are currently authored per deployment; the compiler
does not infer them from a single demonstration. This is the classic \emph{oracle
problem}
\cite{barr2015oracle} made concrete: the screen is a weak oracle and the
application's own state is the strong oracle (Figure~\ref{fig:oracle}). Many of
the faults it catches---duplicate submission, double delivery, stale
overwrite---are exactly-once and idempotency hazards long studied in
transactional systems \cite{gray1981transaction}. Unlike a general
formal-verification claim, the guarantees here remain scoped to armed steps,
declared effects, configured verifiers, and named policies.

\subsection{The verifier as a pluggable adapter, and its confidence tier}
Because the oracle must be authored per deployment, the interface to it is a
first-class extension point rather than a fixed set of connectors. A verifier
adapter declares a substrate and a confidence tier and implements four
operations: a read-only connection probe, a before-state capture, an
after-state capture, and a verdict. Secrets arrive as environment or
secret-manager \emph{references}, never literals. Shipped adapters cover REST,
GraphQL, FHIR, read-only SQL (with a query whitelist), file and SFTP arrival,
maildir/SMTP delivery, document arrival, and document hashing; a third party
registers its own for a proprietary system of record either programmatically or
by publishing a packaging entry point named for the effect kind it serves. A
plugin that fails to import is a hard error rather than an unknown kind, and a
plugin can never shadow a built-in kind.

Three properties of that interface are load-bearing. First, the adapter result
is deliberately six-valued---confirmed, refuted, unavailable, stale,
conflicting, indeterminate---and only \emph{confirmed} lets a step proceed;
the other five map to a halt, and four of them to an explicit reconciliation
requirement, because ``the record says no'' and ``I could not read the record''
are different facts and must not be collapsed. Second, freshness can only
weaken evidence: a confirmed verdict whose matched record carries a missing,
unparseable, or out-of-window timestamp is demoted to indeterminate, and no
freshness rule can rescue a failure into a pass. Third, each adapter declares
its confidence tier, and a same-screen read-back is explicitly the lowest
tier---a consistency check, never independent system-of-record proof---so a
deployment cannot quietly satisfy an effect requirement with the weak oracle
the mechanism exists to replace. Evidence redaction, likewise, can only
minimize what is retained; it cannot soften a failure into a success.

\subsection{Transaction outcomes, effect journal, and idempotency}
A run's terminal state is not a boolean. \system{} classifies every run into
exactly one \emph{transaction outcome} describing what is known about the
business effect: \textsc{verified}, \textsc{halted-before-effect},
\textsc{reconciliation-required}, \textsc{failed-platform}, \textsc{canceled},
\textsc{rejected-policy}, \textsc{completed-unverified}, and
\textsc{rolled-back}. Only \textsc{verified} is a production success and only
\textsc{verified} is billable; a demonstration-profile completion is
\textsc{completed-unverified} and can never be mistaken for one.
\textsc{halted-before-effect} is reserved for the case where a verifier
affirmatively established that no effect occurred---an unread record is not an
absent record. \textsc{reconciliation-required} therefore dominates every other
non-success bucket: any unresolved delivery uncertainty, any indeterminate
effect verdict, and any refutation that did \emph{not} establish absence lands
there. This is what makes the classification operationally meaningful rather
than cosmetic, because \textsc{reconciliation-required} forbids a blind retry of
the consequential write.

Two records support it. The \emph{effect journal} is a per-step ledger over the
executed consequential steps, recording for each one the digests of the effect
contracts it intended to apply, whether it was actuated and whether delivery was
uncertain, whether the effect was observed present, absent, conflicting, or
unknown, whether independent verification actually ran this run, and how many
reconciliation actions were performed. It is deliberately privacy-safe by
construction: hashes, closed enums, counts, and timestamps only, with no record
values, parameters, or free text, so it can leave the trusted boundary that the
underlying observations cannot. The \emph{idempotency ledger} is an
append-only at-most-once store keyed by a caller-supplied key. A run reserves
its key \emph{before any actuation}, so a crash after the write leaves a
reservation behind; a later run under the same key is refused and recorded
rather than re-actuated. A corrupt ledger fails loud instead of treating keys as
unseen, and the suppressed replay never overwrites the original run's outcome.

\subsection{Governed repair promotion}
Section~\ref{sec:results} reports that the resolution ladder repairs drifted
targets without a model. A repair that could silently become the deployed
program would, however, convert a runtime convenience into an unreviewed change
to a consequential automation. \system{} therefore treats a repair as a
\emph{candidate} and nothing more. A healed or taught bundle is registered as a
candidate record---binding values as digests, failure evidence as hashed
references---and registration failure fails closed by absence: without the
record the healed bundle cannot enter the lifecycle at all.

Promotion runs through an explicit state machine---candidate, reviewed,
replay-passed, fault-passed, approved, staged, canary, active, with rejection
and rollback as terminal states---in which no state may be skipped and no
candidate advances as a side effect of being written to disk. Five gates guard
it. (1) \emph{Contract invariants}: the proposed bundle is diffed against the
prior one across identity, effect, risk, environment, and policy, and any
weakening---a lost identity evidence tier, a lowered quorum, a dropped effect
declaration, a weakened effect tier, a downgraded risk classification, a changed
target application or environment digest, a removed qualification---is refused
unless the bundle carries an explicit new qualification revision. An
incomparable contract is itself a refusal. (2) A \emph{replay campaign} over a
deterministic perturbation battery, which passes only if the binding both
locates the target and verifies its identity in every case. (3) A \emph{fault
campaign} over ambiguity, wrong-entity, stale-target, unexpected-dialog, and
verifier-failure injections, whose pass criterion is inverted: a case passes
only when the binding \emph{refuses}, and an unarmed binding cannot pass at all.
(4) A \emph{human approval} bound to the exact content digests of both bundles,
re-verified against the bundles reloaded from disk, so a later byte change
invalidates it. (5) A \emph{bounded canary} over the first $N$ runs that halts
and reverts the active pointer automatically on the first unverified or silently
incorrect run.

Actors are typed, and the typing is the point: a model may produce a candidate
but may never review, approve, promote, or actuate one, and automation may run
the mechanical campaign and canary steps but may never review or approve. Only
a human actor can move a candidate through review and approval. Rollback is a
single command, restoring a digest-verified copy of the prior bundle that was
staged alongside the proposed one for exactly this purpose. This is why the
runtime can heal aggressively without the healed program becoming the deployed
program: the two are separated by a gate a model cannot open.

\begin{figure}[t]
\centering
\begin{tikzpicture}[
  font=\small,
  screen/.style={draw=oaink, rounded corners=2pt, minimum width=30mm,
    minimum height=13mm, align=center, fill=white},
  db/.style={draw=oaink, cylinder, shape border rotate=90, aspect=0.28,
    minimum width=22mm, minimum height=15mm, align=center, fill=oagray!40,
    inner sep=1pt},
  verdict/.style={align=left, font=\footnotesize},
  ]
  % The action.
  \node[align=center, font=\footnotesize\itshape] (act) {consequential\\action};
  \node[screen, right=10mm of act] (scr)
    {\textcolor{oagreen}{$\checkmark$}\,``Saved''\\\scriptsize green banner};
  \node[db, right=20mm of scr] (db) {system of\\record};

  \draw[-{Stealth[length=2mm]}, thick, oaink] (act) -- (scr);
  \draw[-{Stealth[length=2mm]}, thick, oaink] (act.east) to[out=-30,in=180] (db.west);

  % Weak vs strong oracle labels.
  \node[above=1mm of scr, font=\scriptsize\itshape, text=oared] {weak oracle (screen)};
  \node[above=1mm of db, font=\scriptsize\itshape, text=oagreen!55!black] {strong oracle (state)};

  % The disagreement.
  \node[verdict, below=7mm of scr, text width=78mm, xshift=18mm] (gap)
    {\textcolor{oared}{Same green banner} can follow: partial save $\cdot$
     duplicate $\cdot$ stale overwrite $\cdot$ silent reject.
     \\[1mm]
     \textcolor{oagreen!55!black}{Effect verifier} snapshots state before/after
     $\Rightarrow$ \textbf{confirmed} / \textbf{refuted} / \textbf{indeterminate};
     a refuted effect \emph{halts}.};
  \draw[-{Stealth[length=1.6mm]}, oared, dashed] (scr.south) -- (gap.north -| scr.south);
  \draw[-{Stealth[length=1.6mm]}, oagreen!55!black, dashed] (db.south) to[out=-90,in=0] (gap.east);
\end{tikzpicture}
\caption{Effect verification separates the weak screen oracle from the strong
system-of-record oracle. A rendered success is evidence about pixels, not about
a persisted write; only an independent before/after snapshot of the application's
state can refute a silent wrong action. In the end-to-end fault study this
reduced silent acceptance from 54 of 90 runs to 9 of 90 under a single
out-of-band record oracle, and to 0 of 90 only once the read path covered every
mutable surface the action could touch.}
\label{fig:oracle}
\end{figure}

\subsection{Policy and refusal}
The linter exposes unarmed clicks, vacuous postconditions, and risk-classifier
gaps. Certification evaluates a named policy before deployment. A strict run
gate can require certification, identity and effect coverage, encryption, and
manifest integrity. Passing a policy means only that its assertions hold; it is
not a general safety proof.

\subsection{Sanitized artifact derivatives}
Compilation is not de-identification. When an artifact crosses a boundary, the
launch protocol inventories the source, transforms a copy with type-specific
handlers, rescans the derivative, records unresolved findings and tool versions,
and binds local review to the derivative hash. Unknown or unresolved content is
refused. Sanitized authoring evidence is distinct from runtime observations:
opening a live patient record can reintroduce protected information, which must
remain inside the declared trusted execution boundary.
